\documentclass[11pt,a4paper]{article}

\usepackage[T1]{fontenc}
\usepackage[utf8]{inputenc}
\usepackage[french]{babel}
\usepackage{amsmath,amssymb,amsfonts}
\usepackage{graphicx}
\usepackage{booktabs}
\usepackage{multirow}
\usepackage{hyperref}
\usepackage{geometry}
\usepackage{natbib}
\usepackage{enumitem}
\usepackage{float}
\usepackage{caption}
\usepackage{subcaption}
\usepackage{siunitx}

% Inclut la figure si le PNG existe, sinon placeholder (figures gitignorées).
\newcommand{\maybeincludegraphics}[2][]{%
  \IfFileExists{#2}{%
    \includegraphics[#1]{#2}%
  }{%
    \fbox{\parbox[c][5cm][c]{0.85\textwidth}{\centering\textit{Figure à générer}\\ \small\texttt{#2}}}%
  }%
}

\geometry{margin=2.5cm}
\hypersetup{
  colorlinks=true,
  linkcolor=blue,
  citecolor=blue,
  urlcolor=blue
}

\title{%
  \textbf{L'Arène se mord la queue :}\\
  audit temporel, causal et structurel de la vulnérabilité\\
  au style sur \emph{Compar:IA} à travers la loi de Goodhart
}

\author{%
  Équipe hackathon Compar:IA $\times$ Dauphine PSL\\
  \textit{Extension du travail de Simon Zilinskas (2026)}
}

\date{Juin 2026}

\begin{document}
\maketitle

\begin{abstract}
Les arènes d'évaluation de grands modèles de langage (LLM) reposent sur des préférences
humaines pairwise. Lorsqu'un classement devient une cible institutionnelle, la loi de
Goodhart prédit une dérive vers des \emph{proxies} mesurables --- format markdown,
longueur, structure --- au détriment de la qualité intrinsèque.
Nous étendons l'analyse \emph{style-controlled} de \citet{zilinskas2026formatting} sur
\emph{Compar:IA}, arène francophone de la DINUM ($\approx$145\,000 combats nettoyés,
89 modèles), selon cinq axes complémentaires : (i)~dynamique temporelle de la diversité
stylistique (R1) ; (ii)~évolution mensuelle du \emph{style premium} (R2bis) ;
(iii)~preuve causale par réécritures contrefactuelles et juges LLM (R3) ;
(iv)~ablation style \emph{vs.} qualité déclarée et longueur (R5/R5bis) ;
(v)~endogénéité par tier de modèle (R6).
Nos résultats ne confirment \emph{pas} une convergence stylistique ni une dérive
temporelle du premium : la diversité \emph{augmente} ($\beta = +0{,}057$ par mois,
$p < 0{,}001$) et le gras ne monte pas dans le temps ($-1{,}4$ points par mois).
En revanche, la vulnérabilité au formatage est robuste : contrefactuels concis
$\rightarrow$ 89--97\,\% de victoires, verboses $\rightarrow$ 21--27\,\% ;
AUC prédictive du vote passe de 0{,}821 (qualité seule) à 0{,}862 (+ style) ;
après contrôle qualité+longueur, le gras conserve +14{,}5\,\% d'odds de victoire.
Les modèles les mieux classés formatent davantage ($r = 0{,}57$ entre rating et score
de format), excluant une simple stratégie de \emph{gaming} par les modèles faibles.
\emph{Compar:IA} ne présente pas encore les symptômes d'effondrement d'un benchmark
Goodhart, mais reste sensible au style --- un risque à anticiper pour toute métrique
devenue cible.
\end{abstract}

\noindent\textbf{Mots-clés :} évaluation LLM, arène pairwise, loi de Goodhart, style
control, Bradley-Terry, inférence causale, Compar:IA.

\section{Introduction}
\label{sec:intro}

\subsection{Contexte et enjeu}

\emph{Compar:IA} est une plateforme d'évaluation de LLM opérée par la Direction
interministérielle du numérique (DINUM) et le ministère de la Culture
\citep{termignon2026comparia}. Deux modèles anonymes répondent au même prompt ;
l'utilisateur choisit un gagnant ou déclare une égalité. Le classement agrégé
informe la politique publique francophone en matière d'IA.

Ce mécanisme rejoint la famille des arènes pairwise popularisées par Chatbot Arena
\citep{chiang2024chatbot,zheng2023judging}. Le signal utilisateur mêle toutefois
qualité perçue et heuristiques de surface : longueur, markdown, listes, gras
\citep{zheng2023judging,li2024crowdsourced}.

\subsection{Cadre Goodhart}

La loi de Goodhart \citep{goodhart1975problems,manheim2018categorizing} stipule
qu'une mesure utilisée comme cible cesse d'être une bonne mesure. Appliquée aux
benchmarks, elle prévoit trois stades observables \citep{espeland2016engines} :
(i)~convergence des participants vers les traits corrélés au score ;
(ii)~hausse du poids des proxies dans le score ;
(iii)~perte de pouvoir discriminant du benchmark.

Notre question centrale : \emph{Compar:IA} présente-t-il déjà ces signaux, et dans
quelle mesure le formatage influence-t-il les votes indépendamment de la qualité ?

\subsection{Positionnement par rapport à Zilinskas (2026)}

\citet{zilinskas2026formatting} appliquent pour la première fois le \emph{style
control} à une arène non anglophone. Sur un instantané de 145\,096 combats, ils
estiment un modèle de Bradley-Terry \citep{bradley1952rank} où chaque feature de
formatage entre linéairement :
\begin{equation}
  \mathbb{P}(\text{A bat B}) = \sigma\!\left(\beta_A - \beta_B + \sum_f \gamma_f \,\Delta f_{AB}\right),
  \label{eq:bt}
\end{equation}
avec $\Delta f_{AB} = f_A - f_B$ standardisé. Ils trouvent des premiums de
+16--19\,\% par écart-type pour le gras, les listes et les en-têtes ; 76 modèles
changent significativement de rang après correction FDR \citep{benjamini1995controlling},
mais la corrélation globale des classements reste élevée ($r = 0{,}976$).

Nous \textbf{prolongeons} ce travail sans le contredire, en ajoutant :
\begin{enumerate}[noitemsep]
  \item une dimension \textbf{temporelle} (R1, R2bis) absente du papier original ;
  \item une dimension \textbf{causale} par contrefactuels (R3) ;
  \item une décomposition \textbf{qualité / longueur / structure} (R5, R5bis) ;
  \item une analyse d'\textbf{endogénéité} par tier de performance (R6).
\end{enumerate}

\subsection{Thèse empirique retenue}

L'hypothèse initiale du hackathon --- effondrement imminent et convergence
stylistique --- n'est \emph{pas} confirmée par les données. La thèse retenue est
plus nuancée :

\begin{quote}
\emph{Compar:IA} ne montre pas encore les symptômes d'un benchmark en
effondrement Goodhart, mais reste \textbf{vulnérable au style} --- observationnellement
(R5/R5bis/R6) et causalement (R3).
\end{quote}

Le module prospectif R4 (forecast d'effondrement) a été abandonné faute de tendance
à projeter.

\subsection{Plan de l'article}

La Section~\ref{sec:data} décrit les données. La Section~\ref{sec:methods} détaille
les méthodes et calculs. La Section~\ref{sec:results} présente les résultats module
par module. Les Sections~\ref{sec:discussion}--\ref{sec:conclusion} discutent les
implications, limites et conclusions.

\section{Données}
\label{sec:data}

\subsection{Sources}

Nous réutilisons le corpus nettoyé de \citet{zilinskas2026formatting}, publié sur
HuggingFace (\texttt{ministere-culture/comparia-votes}). Après filtrage (paires
identiques, doublons, combats sans texte visible pour modèles \emph{reasoning}),
il reste \textbf{145\,096 combats} impliquant \textbf{89 modèles} avec $\geq 100$
combats chacun, sur 16 mois (octobre 2024 -- janvier 2026).

\subsection{Variables observées}

Pour chaque combat entre modèles $A$ et $B$ sur un prompt $p$ :
\begin{itemize}[noitemsep]
  \item \textbf{Outcome} $y \in \{A, B, \text{tie}\}$ ;
  \item \textbf{Textes} de réponse $r_A, r_B$ ;
  \item \textbf{Métadonnées} : date, mode d'arène (\emph{random}, \emph{custom}, etc.) ;
  \item \textbf{Attributs qualité déclarés} (sous-ensemble) : clarté, exactitude,
        utilité, etc., annotés post-vote par l'utilisateur.
\end{itemize}

\subsection{Features stylistiques}

Cinq features de surface sont extraites par expressions régulières sur $r$ :

\begin{table}[H]
\centering
\caption{Features stylistiques (identiques à Zilinskas).}
\label{tab:features}
\begin{tabular}{lll}
\toprule
Feature & Description & Exemple \\
\midrule
\texttt{headers} & lignes commençant par \# & \texttt{\# Titre} \\
\texttt{lists} & puces ou numérotations & \texttt{- item}, \texttt{1. item} \\
\texttt{bold} & gras markdown & \texttt{**texte**} \\
\texttt{code\_blocks} & blocs de code & \texttt{```python ... ```} \\
\texttt{emoji} & caractères emoji Unicode & \texttt{:smile:} \\
\bottomrule
\end{tabular}
\end{table}

Pour chaque feature $f$, on calcule un compte brut $f(r)$ puis, par combat,
le différentiel $\Delta f = f(r_A) - f(r_B)$. Les versions standardisées
$\Delta f_z$ utilisent la moyenne et l'écart-type empiriques sur l'échantillon.

\subsection{Features qualité et longueur (R5/R5bis)}

Pour R5, nous agrégeons les attributs qualité déclarés en un score
\texttt{quality\_total} (somme des indicateurs binaires positifs par réponse).
Pour R5bis, nous ajoutons $\Delta \log(1 + |r|)$ où $|r|$ est la longueur en
caractères de la réponse assistant (proxy de verbosité, distinct des tokens HF).

\subsection{Échantillons par module}

\begin{table}[H]
\centering
\caption{Tailles d'échantillon effectives par module.}
\label{tab:samples}
\begin{tabular}{lrl}
\toprule
Module & $N$ & Description \\
\midrule
R1 / R1b & 145\,096 & diversité mensuelle, tous combats \\
R2bis & 145\,096 & BT mensuel style-controlled \\
R3 & 100 / 75 & contrefactuels (Mistral seul / ensemble 2 juges) \\
R5 & 79\,073 & combats avec labels qualité \\
R5bis & 79\,073 & idem + longueur \\
R6 & 89 modèles & tiers depuis ratings Zilinskas \\
\bottomrule
\end{tabular}
\end{table}

\section{Méthodes}
\label{sec:methods}

\subsection{R1 --- Diversité stylistique temporelle}

\subsubsection{Indicateur de diversité}

Pour chaque mois $t$, on calcule le centroïde stylistique de chaque modèle $i$ :
$\bar{\mathbf{f}}_{i,t} = \mathbb{E}[f(r) \mid \text{modèle } i, \text{ mois } t]$.
La diversité du mois est la distance moyenne pairwise entre centroïdes actifs :
\begin{equation}
  D_t = \frac{2}{M_t(M_t-1)} \sum_{i<j} \|\bar{\mathbf{f}}_{i,t} - \bar{\mathbf{f}}_{j,t}\|_2.
\end{equation}

Une \textbf{baisse} de $D_t$ signalerait une convergence stylistique (symptôme
Goodhart de stade~1). Nous testons aussi la distance de Wasserstein-2 entre
distributions mensuelles de vecteurs stylistiques.

\subsubsection{Régression temporelle}

Régression OLS :
\begin{equation}
  D_t = \alpha + \beta \cdot \text{month\_idx}_t + \varepsilon_t,
\end{equation}
où $\text{month\_idx}$ est l'index entier du mois depuis l'origine.
$\beta < 0$ compatible avec convergence ; $\beta > 0$ avec diversification.

\subsubsection{Robustesse (R1b)}

Contrôles additionnels : nombre de modèles actifs par mois, volume de combats,
effets fixes de mode d'arène.

\subsection{R2bis --- Style premium longitudinal}

\subsubsection{Modèle BT mensuel}

Chaque mois $t$, on estime \eqref{eq:bt} par maximum de vraisemblance sur les
combats du mois (ties exclus ou modélisés selon variante). Le coefficient
$\gamma_{f,t}$ mesure le \emph{style premium} de la feature $f$ au mois $t$.

\subsubsection{Interprétation en odds}

\begin{equation}
  \text{Premium}_f(\%) = \left(e^{\gamma_f} - 1\right) \times 100.
  \label{eq:premium}
\end{equation}

Un premium de +20\,\% signifie : si $A$ a un écart-type de gras de plus que $B$,
ses odds de victoire augmentent d'environ 20\,\%.

\subsubsection{Tendance temporelle}

Régression de $\gamma_{f,t}$ (ou $\text{Premium}_f$) sur $\text{month\_idx}$ pour
tester si le marché \emph{récompense de plus en plus} le formatage (stade~2 Goodhart).

\subsection{R3 --- Contrefactuels causaux}

\subsubsection{Protocole}

Pour isoler l'effet \emph{causal} du formatage, nous sélectionnons des réponses
neutres (peu de markdown), puis les réécrivons selon deux traitements :
\begin{itemize}[noitemsep]
  \item \textbf{Concise\_direct} : réponse courte, sans structure markdown ;
  \item \textbf{Verbose\_markdown} : réponse longue, listes, gras, en-têtes.
\end{itemize}
Le contenu sémantique est préservé au maximum ; seul le format change.

\subsubsection{Juges LLM}

Deux juges évaluent chaque triplet (prompt, réponse A, réponse B) :
\begin{itemize}[noitemsep]
  \item \textbf{Mistral Large} (API Mistral) ;
  \item \textbf{Llama-3.3-70B} (API Groq).
\end{itemize}

\subsubsection{Filtre de similarité sémantique}

Pour éviter que le juge tranche sur un changement de contenu, nous exigeons
$\cos(\text{emb}_A, \text{emb}_B) \geq 0{,}85$ (embeddings sentence-transformers).

\subsubsection{Statistiques}

Taux de victoire du traitement avec intervalle de confiance de Wilson
\citep{wilson1927probable} à 95\,\%. Accord inter-juges : Cohen's $\kappa$
\citep{cohen1960coefficient} et taux de victoire sur les paires où les deux juges
sont valides et d'accord sur la comparaison.

L'effet causal estimé (ATE) pour le traitement $T \in \{\text{concise}, \text{verbose}\}$ :
\begin{equation}
  \widehat{\text{ATE}}_T = \frac{1}{N}\sum_{i=1}^{N} \mathbb{1}[\text{juge préfère } T^{(i)}].
\end{equation}

\subsection{R5 --- Style \emph{vs.} qualité déclarée}

\subsubsection{Modèle logistique}

Sur les combats avec attributs qualité, outcome binaire $y_i = 1$ si A gagne :
\begin{equation}
  \log\frac{p_i}{1-p_i} = \beta_0 + \beta_q Q_i + \sum_f \beta_f \Delta f_{i} + \varepsilon_i,
  \label{eq:logit}
\end{equation}
où $Q_i$ est le différentiel de score qualité agrégé.

\subsubsection{Ablations et AUC}

Nous comparons l'AUC in-sample de plusieurs spécifications :
qualité seule, style seul, qualité+style, avec interactions temporelles.
L'\textbf{incrément} $\Delta\text{AUC}$ mesure le gain prédictif du style \emph{après}
contrôle de la qualité déclarée.

\subsubsection{Coefficients marginaux}

Les odds ratios par feature :
$\text{OR}_f = e^{\beta_f}$, interprétés comme en \eqref{eq:premium}.

\subsection{R5bis --- Structure \emph{vs.} longueur}

R5bis étend \eqref{eq:logit} :
\begin{equation}
  \log\frac{p_i}{1-p_i} = \beta_0 + \beta_q Q_i + \beta_l \Delta\log(1+|r|)_i + \sum_f \beta_f \Delta f_i.
  \label{eq:logit-r5bis}
\end{equation}

Objectif : vérifier si le premium de structure (gras, listes\ldots) persiste une
fois la longueur contrôlée --- réconciliant R3 (effet massif du format) et R5
(qualité domine en AUC).

\subsection{R6 --- Endogénéité par tier}

\subsubsection{Construction des tiers}

Les modèles sont stratifiés en quartiles selon le rating Bradley-Terry de
Zilinskas (JSON de ratings style-controlled).

\subsubsection{Statistiques}

\begin{itemize}[noitemsep]
  \item Corrélation de Pearson $r$ entre rating et score de format moyen ;
  \item Différence de premium gras entre top et bottom tier ;
  \item Régression avec interaction $\gamma_f \times \text{tier}$.
\end{itemize}

Si seuls les modèles faibles \emph{gameaient} le formatage, on s'attendrait à un
premium concentré en bas de classement. L'hypothèse alternative (modèles forts
mieux formatés \emph{et} sur-récompensés) prédit $r > 0$ et premium croissant avec le tier.

\section{Résultats}
\label{sec:results}

\subsection{R1 --- Pas de convergence stylistique}

La diversité $D_t$ \textbf{augmente} dans le temps (Figure~\ref{fig:r1}) :
\begin{equation}
  \hat{\beta} = +0{,}057 \text{ par mois}, \quad p < 0{,}001.
\end{equation}

Après contrôles (R1b), le signe positif persiste. La distance de Wasserstein entre
mois consécutifs ne montre pas de contraction systématique.

\begin{figure}[H]
  \centering
  \maybeincludegraphics[width=0.85\textwidth]{figures/R1_convergence_robustness.png}
  \caption{Diversité stylistique mensuelle (R1). Tendance positive : pas de convergence
  vers un format unique.}
  \label{fig:r1}
\end{figure}

\textbf{Interprétation.} Le stade~1 Goodhart (convergence) n'est pas observé ; les
modèles restent stylistiquement distincts, voire plus dispersés.

\subsection{R2bis --- Premium élevé mais stable dans le temps}

En niveau, le style premium est \textbf{fort} et cohérent avec Zilinskas : le gras
vaut environ +20\,\% d'odds par écart-type (Figure~\ref{fig:r2bis}).

En revanche, la pente temporelle du premium gras est \textbf{négative} :
$\hat{\beta} = -1{,}4$ points de premium par mois --- le formatage ne gagne \emph{pas}
en influence relative au fil des mois.

\begin{figure}[H]
  \centering
  \maybeincludegraphics[width=0.85\textwidth]{figures/R2bis_style_premium_longitudinal.png}
  \caption{Évolution mensuelle du style premium (R2bis). Niveau élevé, pente non positive.}
  \label{fig:r2bis}
\end{figure}

\textbf{Interprétation.} Vulnérabilité structurelle au style (niveau), sans dérive
temporelle du type stade~2 Goodhart.

\subsection{R3 --- Preuve causale}

Table~\ref{tab:r3} et Figure~\ref{fig:r3} résument les contrefactuels.

\begin{table}[H]
\centering
\caption{Résultats contrefactuels (R3). IC 95\,\% Wilson.}
\label{tab:r3}
\begin{tabular}{lrrrrr}
\toprule
Comparaison & $N$ & Mistral & Llama-70B & Accord 2/2 & IC accord \\
\midrule
Concise vs neutre & 45 & 88{,}9\,\% & 90{,}9\,\% & 97{,}3\,\% & [86, 100] \\
Verbose vs neutre & 30 & 26{,}7\,\% & 33{,}3\,\% & 26{,}9\,\% & [14, 46] \\
\bottomrule
\end{tabular}
\end{table}

\begin{figure}[H]
  \centering
  \maybeincludegraphics[width=0.85\textwidth]{figures/R3_ensemble_judge.png}
  \caption{Taux de victoire par traitement et juge (R3). $\kappa(\text{Mistral}, \text{Llama-70B}) = 0{,}67$.}
  \label{fig:r3}
\end{figure}

Changer \emph{uniquement} le format suffit à inverser massivement les préférences :
concis bat neutre dans $\approx$90--97\,\% des cas ; verbose \emph{perd} contre neutre
($\approx$27\,\%). L'accord inter-juges est substantiel ($\kappa = 0{,}67$).

\textbf{Limites.} Juges LLM, non utilisateurs Compar:IA ; $N$ modeste mais effets
très grands.

\subsection{R5 --- Le style ajoute de l'AUC au-delà de la qualité}

Table~\ref{tab:r5auc} : ablation AUC in-sample sur 79\,073 combats.

\begin{table}[H]
\centering
\caption{AUC in-sample par spécification logistique (R5).}
\label{tab:r5auc}
\begin{tabular}{lrr}
\toprule
Spécification & AUC & $d$ features \\
\midrule
Style seul (agrégé) & 0{,}667 & 1 \\
Style seul (composantes) & 0{,}669 & 5 \\
Qualité seule (agrégé) & 0{,}811 & 1 \\
Qualité seule (composantes) & 0{,}821 & 7 \\
Qualité + style (agrégé) & 0{,}858 & 2 \\
Qualité + style (composantes) & 0{,}862 & 12 \\
Qualité + style + interactions mois & 0{,}863 & 17 \\
\bottomrule
\end{tabular}
\end{table}

La qualité déclarée domine ($\text{AUC}_q = 0{,}821$), mais le style apporte
$+0{,}041$ AUC. Après contrôle qualité, le gras conserve $\approx +30{,}6\,\%$
d'odds de victoire (Figure~\ref{fig:r5}).

\begin{figure}[H]
  \centering
  \maybeincludegraphics[width=0.85\textwidth]{figures/R5_style_vs_quality.png}
  \caption{Coefficients style marginaux après contrôle qualité (R5).}
  \label{fig:r5}
\end{figure}

\subsection{R5bis --- Structure vs longueur}

Une fois qualité \emph{et} $\Delta\log(1+|r|)$ contrôlés (Table~\ref{tab:r5bis}) :

\begin{table}[H]
\centering
\caption{R5bis : ablation structure vs longueur ($N = 79\,073$).}
\label{tab:r5bis}
\begin{tabular}{lrr}
\toprule
Spécification & AUC & Premium gras résiduel \\
\midrule
Qualité + longueur & 0{,}865 & --- \\
Qualité + longueur + style & 0{,}865 & +14{,}5\,\% \\
\bottomrule
\end{tabular}
\end{table}

La longueur explique une partie de l'effet «~verbose~» de R3 (+40{,}5\,\% odds pour
$\Delta\log(1+|r|)$), mais le gras conserve un premium de +14{,}5\,\%.
R3 et R5 se réconcilient : format \emph{et} longueur opèrent, qualité reste le
prédicteur dominant en AUC.

\begin{figure}[H]
  \centering
  \maybeincludegraphics[width=0.85\textwidth]{figures/R5bis_structure_vs_length.png}
  \caption{Décomposition structure vs longueur (R5bis).}
  \label{fig:r5bis}
\end{figure}

\subsection{R6 --- Les tops formatent plus (pas un gaming des faibles)}

\begin{itemize}[noitemsep]
  \item $r(\text{rating}, \text{format score}) = 0{,}57$ ;
  \item premium gras : bottom tier vs top tier = +41{,}6 points ;
  \item les modèles bien classés produisent \emph{plus} de markdown, et en
        bénéficient davantage.
\end{itemize}

\begin{figure}[H]
  \centering
  \maybeincludegraphics[width=0.85\textwidth]{figures/R6_endogeneity_tiers.png}
  \caption{Style premium et formatage moyen par tier de rating (R6).}
  \label{fig:r6}
\end{figure}

\textbf{Interprétation.} L'endogénéité est \emph{symétrique} : le biais de formatage
n'est pas l'apanage des modèles faibles tentant de compenser ; il touche aussi les
leaders dont le style est à la fois plus riche et mieux récompensé.

\subsection{Synthèse des modules}

\begin{table}[H]
\centering
\caption{Cartographie des résultats et lien avec Goodhart.}
\label{tab:synth}
\small
\begin{tabular}{llll}
\toprule
Module & Question & Résultat & Stade Goodhart \\
\midrule
R1 & Convergence ? & Diversité $\uparrow$ & Stade 1 : non \\
R2bis & Premium $\uparrow$ dans le temps ? & Niveau haut, pente $\leq 0$ & Stade 2 : non \\
R3 & Format cause-t-il le vote ? & Oui (large ATE) & Vulnérabilité \\
R5/R5bis & Au-delà qualité+longueur ? & Oui (+14--31\,\% gras) & Vulnérabilité \\
R6 & Gaming des faibles ? & Non ($r=0{,}57$) & Nuance \\
R4 & Effondrement imminent ? & Abandonné & --- \\
\bottomrule
\end{tabular}
\end{table}

\section{Discussion}
\label{sec:discussion}

\subsection{Ce que nous apprenons sur Compar:IA}

Trois messages pour les responsables de la plateforme :
\begin{enumerate}
  \item \textbf{Pas de panique temporelle} : ni convergence ni inflation récente
        du premium ; le benchmark n'est pas en train de s'effondrer \emph{observablement}.
  \item \textbf{Vulnérabilité réelle} : formatage et longueur influencent les votes
        causalement et prédictivement, même contrôlant la qualité déclarée.
  \item \textbf{Biais structurel} : les modèles en tête de classement sont aussi
        ceux qui formatent le plus ; un simple contrôle style ne suffit pas à
        «~corriger~» le classement sans modèle causal plus fin.
\end{enumerate}

\subsection{Structure vs longueur}

R3 montrait un contraste concis/verbose extrême ; R5 montrait la domination de la
qualité en AUC. R5bis montre que ces résultats ne se contredisent pas : la longueur
est un mediateur majeur, mais la structure markdown conserve un effet propre ---
exactement le pattern rapporté par LMSYS Style Control sur données anglaises
\citep{li2024crowdsourced}.

\subsection{Qualité déclarée comme confondant}

Les labels qualité sont \emph{post-vote} et corrélés à la préférence. R5 contrôle
ce biais observable ; il ne capture pas la qualité non déclarée. L'AUC de 0{,}862
reste in-sample ; une validation croisée temporelle serait nécessaire pour une
 généralisation out-of-sample.

\subsection{Recommandations}

\begin{enumerate}
  \item Publier des classements \emph{style-controlled} aux côtés des classements bruts
        (suivre Zilinskas) ;
  \item Randomiser ou normaliser la longueur dans une fraction des combats ;
  \item Monitorer $D_t$ et $\gamma_{f,t}$ mensuellement comme \emph{early warnings} ;
  \item Compléter les juges LLM (R3) par des contrefactuels annotés par utilisateurs réels.
\end{enumerate}

\section{Limites}
\label{sec:limits}

\begin{itemize}
  \item \textbf{Labels qualité post-hoc} (R5) : biais de cohérence avec le vote.
  \item \textbf{AUC in-sample} : optimisme possible ; pas de split temporel strict.
  \item \textbf{R3 juges LLM} : externalité limitée vers votes humains Compar:IA.
  \item \textbf{Longueur en caractères} (R5bis) : pas en tokens ; corrélation forte
        mais imparfaite avec la facturation API.
  \item \textbf{Tiers R6} : ratings importés du snapshot Zilinskas, non ré-estimés
        mensuellement.
  \item \textbf{Causalité R3} : réécritures automatiques ; filtre cosine $\geq 0{,}85$
        arbitraire.
  \item \textbf{16 mois} : fenêtre courte pour des phénomènes Goodhart lents
        \citep{espeland2016engines}.
\end{itemize}

\section{Conclusion}
\label{sec:conclusion}

Ce travail prolonge l'audit de \citet{zilinskas2026formatting} sur \emph{Compar:IA}
selon une grille temporelle, causale et structurelle inspirée de la loi de Goodhart.
Les données \textbf{ne montrent pas} une convergence stylistique ni une dérive
temporelle du style premium ; elles \textbf{montrent} une sensibilité persistante
au formatage --- causalement (R3), prédictivement (R5/R5bis) et corrélée au rang
(R6).

Pour la gouvernance des benchmarks LLM, la leçon est proactive : un classement peut
rester stable en corrélation ($r = 0{,}976$ chez Zilinskas) tout en restant
gameable sur la marge. \emph{Compar:IA} n'a pas encore «~mordu sa queue~», mais
l'arène récompense déjà des proxies de surface --- il est temps de construire des
garde-fous avant qu'une cible explicite ne transforme ces vulnérabilités en
stratégie d'optimisation.

\section*{Remerciements}

Nous remercions Simon Zilinskas pour son travail préalable, le partage de code et
de données, ainsi que l'équipe Compar:IA/DINUM pour la mise à disposition des jeux
de données ouverts.

\bibliographystyle{plainnat}
\bibliography{refs}

\appendix

\section{Formules complémentaires}
\label{app:formulas}

\subsection{Intervalle de Wilson}

Pour $k$ succès sur $n$ essais, avec $z = 1{,}96$ :
\begin{align}
  \tilde{p} &= \frac{k + z^2/2}{n + z^2}, \\
  \text{IC}_{95\%} &= \tilde{p} \pm \frac{z}{n + z^2}\sqrt{\frac{k(n-k)}{n} + \frac{z^2}{4}}.
\end{align}

\subsection{Cohen's $\kappa$}

\begin{equation}
  \kappa = \frac{p_o - p_e}{1 - p_e},
\end{equation}
où $p_o$ est l'accord observé et $p_e$ l'accord attendu par indépendance.

\subsection{Distance de Wasserstein (R1)}

Pour deux distributions empiriques $\hat{\mu}_t, \hat{\mu}_{t+1}$ sur les vecteurs
stylistiques :
\begin{equation}
  W_2(\hat{\mu}_t, \hat{\mu}_{t+1}) = \left(\inf_{\gamma \in \Gamma} \mathbb{E}_{(x,y)\sim\gamma}\|x-y\|_2^2\right)^{1/2}.
\end{equation}

\section{Reproduction des résultats}
\label{app:repro}

Scripts principaux (répertoire \texttt{scripts/}) :
\begin{itemize}[noitemsep]
  \item \texttt{r1\_robustness.py}, \texttt{r2bis\_style\_premium.py}
  \item \texttt{r3\_ensemble\_judge.py}, \texttt{r5\_style\_vs\_quality.py}
  \item \texttt{r5bis\_structure\_vs\_length.py}, \texttt{r6\_endogeneity\_tiers.py}
\end{itemize}

Données : \texttt{data/raw/hf/comparia-votes/votes.parquet}.
Figures : \texttt{paper/figures/}. Tables : \texttt{paper/tables/}.

\end{document}
